% TOP_PROBLEM: {"id": 3173, "title": "Strong 5-cycle double cover conjecture", "category_id": 3, "category": {"id": 3, "name": "graph_theory", "display_name": "Graph Theory", "description": "Problems involving graphs, networks, and their properties.", "slug": "graph-theory", "order_index": 3, "created_at": "2026-07-31T15:26:25.671Z"}, "status": "open", "problem_number": "OPG-37241", "set_id": 12}
% FIRST_POSTED: 2026-10-01
% ATTEMPT_STATUS: unresolved
% UnsolvedMath ID 3173; source catalogue code OPG-37241
% !TeX program = lualatex
% GPT-6 Astra Ultra research attempt; independent partial work, 2026-10-01.
\documentclass[11pt,a4paper]{article}
\usepackage[margin=25mm]{geometry}
\usepackage{fontspec}
\setmainfont{lmroman10-regular.otf}[BoldFont=lmroman10-bold.otf,
 ItalicFont=lmroman10-italic.otf,BoldItalicFont=lmroman10-bolditalic.otf]
\usepackage{amsmath,amssymb,mathtools}
\usepackage{unicode-math}
\setmathfont{latinmodern-math.otf}
\AtBeginDocument{\let\setminus\smallsetminus}
\usepackage{xurl,hyperref}
\hypersetup{colorlinks=true,urlcolor=blue}
\setcounter{secnumdepth}{0}
\setlength{\parindent}{0pt}
\setlength{\parskip}{5pt}
\setlength{\emergencystretch}{3em}
\begin{document}
\section*{Strong 5-cycle double cover conjecture}\label{3173}
{\small UnsolvedMath ID 3173 · OPG-37241 · Graph Theory · OpenGarden}
\subsection{Definitions and mathematical statement}
Let $G=(V,E)$ be a finite simple cubic graph with no bridge,
and fix a circuit $C$ of $G$. Here cubic means degree three
at every vertex, a bridge is an edge whose removal increases
the number of connected components, and a circuit is one
connected simple cycle.

For this problem a covering cycle means an edge set
$F\subseteq E$ for which every vertex of $(V,F)$ has degree
zero or two. Thus its nontrivial components are disjoint
circuits, and $F$ may be empty. A five-cycle double cover
is a list $F_1,\ldots,F_5$ of such edge sets satisfying
\[
 \sum_{i=1}^5\mathbf1_{\{e\in F_i\}}=2
 \qquad(e\in E).
\]
Repetitions and empty members are allowed, so the phrase
``five'' means at most five nonempty covering members.

The strong conjecture requires that for every pair $(G,C)$
one can choose the cover with $E(C)\subseteq F_j$ for some
index $j$. Since the degrees inside $F_j$ are at most two,
the given circuit then occurs as a whole component of
$F_j$; additional disjoint circuits may occur there too.

It is essential that the five members need not themselves
be connected. The statement does not ask for five individual
circuits, and merely ensuring that the edges of $C$ appear
somewhere in the double cover is insufficient. All edges
of the specified circuit must belong to one common member.
\subsection{Short English statement}
Can every prescribed circuit in a bridgeless cubic graph be a component of one member of a five-cycle double cover?
\subsection{Research attempt: an even-factor criterion and a Petersen certificate}
\paragraph{Scope and current literature.}
GPT-6 Astra Ultra, 1 October 2026. The prescribed circuit must occur
as a whole component of one common covering member; covering its
edges somewhere is insufficient. The source problem [S2] records this
stronger condition. Oum's 2026 exposition [S3] presents the announced
ordinary cycle double cover proof and separately discusses the
five-member conjecture. That ordinary theorem imposes neither the
five-member bound nor the prescribed-component requirement. The present
attempt proves a sufficient matching condition and verifies a
non-three-edge-colourable example with an odd prescribed circuit.

\paragraph{1. A useful sufficient condition.}
Suppose $G$ has a perfect matching $M$ disjoint from $E(C)$ such
that every component of the complementary two-factor $F=E(G)\setminus M$
has even length. Since every vertex of $C$ already has its two
$C$-edges in $F$, the circuit $C$ is a whole component of $F$.
Alternately colour each component of $F$ with colours one and two,
and colour all matching edges with colour three. This is a proper
three-edge-colouring: each vertex has one incident edge of each colour.

Let $F_{12},F_{13},F_{23}$ be the unions of the indicated colour
classes. Each has degree two at every vertex, so each is an allowed
covering member. Each edge belongs to exactly two of these three
members. Moreover $F_{12}=F$ contains $C$ as a whole component.
Appending two empty members gives a five-cycle double cover satisfying
the prescribed-circuit condition. This proves the strong target whenever
the stated matching and even-factor condition can be certified.

In particular, if $C$ is Hamiltonian, the unused edges of the cubic
graph form a perfect matching. A cubic graph has an even number
of vertices, by the degree sum $3|V|=2|E|$, so this Hamilton cycle
has even length. The condition applies and proves the target for
every prescribed Hamilton circuit. For a nonspanning even circuit,
its parity alone does not guarantee the required matching or the
parities of the other components.

\paragraph{2. An explicit odd-circuit example.}
Use Petersen vertices $0,\ldots,9$ and index its edges as follows:
\[
\begin{array}{c|rrrrrrrr}
 j&0&1&2&3&4&5&6&7\\
 e_j&01&12&23&34&04&05&16&27\\[2pt]
 j&8&9&10&11&12&13&14&\\
 e_j&38&49&57&68&79&58&69&
\end{array}
\]
Here $ab$ denotes the unordered edge with endpoints $a,b$.
Prescribe the outer five-cycle $C$ with edge indices $0,1,2,3,4$.
The following five sets are a certificate:
\[
\begin{aligned}
 F_1&=\{0,1,2,3,4\},\\
 F_2&=\{0,2,5,6,7,8,10,11\},\\
 F_3&=\{3,8,9,10,12,13\},\\
 F_4&=\{1,6,7,12,14\},\\
 F_5&=\{4,5,9,11,13,14\}.
\end{aligned}
\]
Interpreting each index as its listed edge, every vertex has degree
zero or two in each $F_i$. Counting the occurrences of every index
$0,\ldots,14$ gives exactly two, and $F_1$ is precisely the
prescribed circuit. The script [S4] checks these statements directly,
and also enumerates all even edge sets and reproduces a search for
such a certificate. The example demonstrates that the three-colouring
criterion is sufficient but not necessary.

\paragraph{3. A genuine limit of the colouring approach.}
A component using only two colours in a proper edge-colouring must
alternate them and therefore have even length. An odd prescribed
circuit can never be a component of any of the three bicoloured
members produced in paragraph 1. Furthermore the same script checks
that all six perfect matchings of Petersen have complements consisting
of two five-cycles, ruling out a proper three-edge-colouring there.
The explicit five-member certificate overcomes both limitations on this
one input, without supplying a general construction for odd circuits.

\paragraph{Verification and first remaining gap.}
The edge list fixes all labelling conventions and the certificate
contains no implicit graphical choices. The search is finite and
uses exact parity and incidence counts; it establishes only this
specified Petersen pair $(G,C)$. For an arbitrary prescribed circuit,
no matching with the required even complement is guaranteed, and
no argument here bounds a more general repair procedure by five
members while keeping $C$ intact.

\subsection{Final status}
UNRESOLVED. The even-factor criterion, all prescribed Hamilton circuits,
and the displayed odd Petersen circuit are handled. The universal
strong five-cycle double cover conjecture is not proved. Completion
estimate: no defensible global percentage.

\subsection{Sources}
\begin{itemize}
\item[S1] ULAM.AI and the UnsolvedMath Contributors, supplied problems.json, record 3173 (OPG-37241), OpenGarden collection. Catalogue identity and source transcription; CC BY 4.0.
\url{https://huggingface.co/datasets/ulamai/UnsolvedMath}.
\item[S2] Open Problem Garden, Strong 5-cycle double cover conjecture. Original problem and discussion; the mathematical formulation below is expanded from the supplied source record.
\url{http://www.openproblemgarden.org/op/strong_5_cycle_double_cover_conjecture}.
\item[S3] S.-i. Oum, \emph{A proof of the cycle double cover conjecture by OpenAI: An exposition}, arXiv:2607.16356v3 (2026), especially Section 9.2. Checked 1 October 2026.
\url{https://arxiv.org/abs/2607.16356}.
\item[S4] GPT-6 Astra Ultra, exact Petersen certificate and exhaustive finite search, 1 October 2026: \texttt{scripts/check\_attempt\_batch02\_connectivity\_2026\_10\_01.py} in this repository; run with Python 3.
\end{itemize}
\end{document}
